\documentclass[11pt]{amsart}
\usepackage[T1]{fontenc}
\usepackage{newtxtext,newtxmath}
\usepackage[margin=1in]{geometry}
\raggedbottom
\usepackage{booktabs,longtable,fancyvrb}
\usepackage[hidelinks]{hyperref}
\newtheorem{theorem}{Theorem}[section]
\newtheorem{lemma}[theorem]{Lemma}
\theoremstyle{definition}
\newtheorem{definition}[theorem]{Definition}
\title[Consecutive primes at distance 122,290]{Consecutive primes at distance 122,290}
\author{Will Howe}
\email{howew@purdue.edu}
\author{Grant Huffman}
\email{huffma97@purdue.edu}
\author{Jensen Kohlmeyer}
\email{jenwin@purdue.edu}
\author{Liam Kruer}
\email{kruerl@purdue.edu}
\date{November 16, 2026}
\subjclass[2020]{Primary 11A41; Secondary 11Y11, 11Y55}
\keywords{Prime gaps, covering congruences, primality certificates}
\begin{document}
\begin{abstract}
We exhibit two consecutive primes whose difference is exactly $122\,290$.
Both primes have $8\,262$ decimal digits. Every integer strictly between them
has a prime divisor at most $32\,467$. We give the covering argument and the
elliptic-curve criterion that certify the result. The proof is computer-assisted and uses finite arithmetic checks to establish
endpoint primality and the compositeness of every interior integer.
\end{abstract}
\maketitle

\section{The result and the covering argument}
A prime gap of length $h$ is a pair of primes $p,p+h$ such that every integer
between them is composite. We prove the following statement.

\begin{theorem}\label{thm:main}
Let $p$ be the integer written in Appendix~\ref{app:p}. Then $p$ and
$p+122\,290$ are consecutive primes, each with $8\,262$ decimal digits.
Moreover, every integer $n$ with $p<n<p+122\,290$ has a prime divisor at
most $32\,467$.
\end{theorem}

The full residue table and the two primality chains form the embedded
arithmetic appendix to this PDF. Their mathematical format is specified below.

A covering congruence supplies one small divisor for each interior integer.

\begin{lemma}\label{lem:cover}
Let $h\ge2$, let $\mathcal Q$ be a finite set of primes, and choose a residue
$b_q$ modulo each $q\in\mathcal Q$. Put $W=\prod_{q\in\mathcal Q}q$, and
let $0\le a<W$ satisfy $a\equiv b_q\pmod q$ for every $q\in\mathcal Q$.
Suppose that
\[
 \{1,\ldots,h-1\}\subseteq
 \bigcup_{q\in\mathcal Q}\{j:j\equiv-b_q\pmod q\}.
\]
If $m\ge1$, $p=a+mW>\max\mathcal Q$, and $p,p+h$ are prime, then they are
consecutive primes.
\end{lemma}
\begin{proof}
For $1\le j<h$, choose $q\in\mathcal Q$ with $j\equiv-b_q\pmod q$.
Since $p\equiv a\equiv b_q\pmod q$, we have $q\mid p+j$.
The inequalities $1<q<p+j$ show that $p+j$ is composite.
\end{proof}

For the integer $p$ in Theorem~\ref{thm:main}, the residue table has
$2\,180$ pairs $(q,b_q)$ and satisfies
\[
 b_q\equiv p\pmod q,\qquad b_q\not\equiv0,-122\,290\pmod q.
\]
Every $q$ is prime and $q\le32\,467$. The Chinese remainder construction is
$p=a+530\,520W$, where $W=\prod q$, $0\le a<W$, and
$a\equiv b_q\pmod q$. The integer $W$ has $8\,256$ decimal digits.

To check the interior directly, start with $S=\{1,\ldots,122\,289\}$. For each prime
$q\le B$, remove the offsets $j\equiv-p\pmod q$. The successive numbers
of remaining offsets are
\begin{center}
\begin{tabular}{r r@{\qquad}r r}
\toprule
$B$ & Remaining & $B$ & Remaining\\
\midrule
$2$ & $61\,144$ & $14\,983$ & $818$\\
$97$ & $14\,577$ & $19\,997$ & $126$\\
$997$ & $8\,665$ & $24\,989$ & $26$\\
$4\,999$ & $4\,673$ & $29\,989$ & $3$\\
$9\,973$ & $2\,342$ & $32\,467$ & $0$\\
\bottomrule
\end{tabular}
\end{center}
Thus every interior integer has a proper prime divisor at most $32\,467$.

\section{Endpoint primality and completion of the proof}
We use the following form of the elliptic-curve primality criterion
\cite{atkinmorain}. All group operations in the statement are performed
modulo~$N$.

\begin{lemma}\label{lem:ecpp}
Let $N>3$ be coprime to $6$, and let
$E:y^2=x^3+Ax+B$ satisfy $\gcd(4A^3+27B^2,N)=1$.
Suppose that $r$ is prime, $r>(N^{1/4}+1)^2$, and an affine point $R$ on $E$
satisfies $[r]R=O$. If the computation uses only denominators coprime to $N$,
then $N$ is prime.
\end{lemma}
\begin{proof}
If $N$ were composite, it would have a prime divisor $\ell\le\sqrt N$.
Reduction modulo $\ell$ gives a nonsingular elliptic curve, and $R$ reduces
to an affine point. This point has order $r$, since it is nonzero and is
annihilated by the prime $r$. Hasse's bound therefore gives
\[
 r\le\#E(\mathbb F_\ell)
 \le(\sqrt\ell+1)^2
 \le(N^{1/4}+1)^2,
\]
a contradiction.
\end{proof}

For each endpoint, the certificate is a finite chain of integers $N_i$ and
tuples $(t_i,s_i,A_i,X_i,Y_i)$, with
\[
 N_{i+1}=\frac{N_i+1-t_i}{s_i},\qquad
 B_i\equiv Y_i^2-X_i^3-A_iX_i\pmod{N_i}.
\]
For $0\le i<k$, $N_i>3$, $s_i>0$, and $N_{i+1}$ is an integer.
At each step we have
\begin{equation}\label{eq:checks}
\begin{gathered}
 \gcd(N_i,6)=\gcd(4A_i^3+27B_i^2,N_i)=1,\\
 N_{i+1}>(N_i^{1/4}+1)^2,\\
 R_i=[s_i](X_i,Y_i)\ne O,\qquad [N_{i+1}]R_i=O.
\end{gathered}
\end{equation}
The point $R_i$ is affine and every division in its computation, and in that
of $[N_{i+1}]R_i$, has denominator coprime to $N_i$. These are exact modular
arithmetic checks. The two chains have the following endpoints.
\begin{center}
\begin{tabular}{l r r}
\toprule
Starting integer & Number of steps & Terminal integer\\
\midrule
$p$ & $647$ & $77\,611\,968\,613$\\
$p+122\,290$ & $659$ & $78\,278\,019\,539$\\
\bottomrule
\end{tabular}
\end{center}
Both terminal integers are prime. For a direct order check, their
factorizations and bases are
\[
\begin{array}{rcll}
77\,611\,968\,613-1&=&2^2\cdot3^4\cdot937\cdot255\,649,&g=6,\\
78\,278\,019\,539-1&=&2\cdot389\cdot5\,531\cdot18\,191,&g=2.
\end{array}
\]
Every displayed factor is prime by trial division. In each row,
$g^{N-1}\equiv1\pmod N$ and
$\gcd(g^{(N-1)/r}-1,N)=1$ for every prime $r\mid N-1$.
Consequently $g$ has order $N-1$ modulo every prime divisor $\ell$ of $N$,
so $N-1\mid\ell-1$ and $\ell\ge N$. This proves the claim.

\begin{proof}[Proof of Theorem~\ref{thm:main}]
The finite modular checks \eqref{eq:checks}, starting from the two terminal
primes and applying Lemma~\ref{lem:ecpp} backwards along each chain, prove
that $p$ and $p+122\,290$ are prime. The covering calculation in Section~1
proves every intermediate integer composite. Thus the primes are consecutive.
The decimal representation of $p$ and addition of $122\,290$ give $8\,262$
digits in both cases.
\end{proof}

\subsection*{Acknowledgment}
The authors acknowledge OpenAI's AI systems for assistance with the
construction, finite arithmetic verification, and preparation of this paper.

\begin{thebibliography}{1}
\bibitem{atkinmorain}
A.~O.~L. Atkin and F. Morain,
\emph{Elliptic curves and primality proving},
Math. Comp. \textbf{61} (1993), no.~203, 29--68.
\end{thebibliography}

\appendix
\section{The starting prime}\label{app:p}
The decimal expansion of $p$ is displayed below. The lines are concatenated
and spaces separate blocks of ten digits. The integer $p+122\,290$ is then
determined by addition.
\begin{Verbatim}[fontsize=\fontsize{6.4}{7.3}\selectfont]
2524630355 5392095200 6849031818 2436336455 8705070412 1679587043 2398027747 9814213567 5465079894 0493256636 4749901363 7823525575
4209929759 9738020881 7752842148 3859541993 4209196681 4556241831 2586384700 5663906315 1839543596 7260996320 6698569079 2712486548
1453992286 7309146566 4147352403 1034075473 0932486503 7771948675 0538956940 2856290174 5774135130 8828043496 4853408946 0162269581
4061115880 3828627806 0403256665 7074998453 7470356071 3668275390 8061279014 2506972619 0119303941 2589965001 5584117686 9475615343
6449533093 2910451236 1220951615 2184713873 2797330447 6485527212 2577310486 5208352493 4033071492 6223893449 5446224729 9128489789
0064275766 6889877248 3321933151 8539283669 5896459721 0204974914 1353610729 4232034593 9209438523 9061891081 3827634985 7887852391
0418366629 1542255791 1958492222 4598665039 5249055786 2846777088 8395643767 5937390640 5292674126 2373670525 4827389933 0732408098
6939913760 1274338263 7517970302 1828395992 8627815024 5649001329 3631552984 9105155093 4575288391 0177794671 8860046765 0792351466
3655060518 8431946824 8646326914 5199183796 4304679065 2821823293 0006079496 1699777199 0824081861 5185650275 0386743802 5538811771
5051917928 6453527121 4754385031 2847869816 3076786496 1782064058 3458729594 5550585916 0437474672 4563880450 4628050102 8610642121
9866783775 7677493203 4046536681 6758671044 9913134478 4319417191 2897639301 3038574757 3431939458 9607257413 4622799655 1751800158
4631126739 7168658370 9542877343 5317690685 0814908329 5618131314 3676935022 7382996028 3689035536 6199596159 8289943849 7166392882
1564506423 4457934027 8485098886 4962041472 1612544564 9886958936 1334402470 7084549558 7622917919 2961450426 9004069963 9871721950
5722466388 2694631561 4433503669 4076733685 8200918093 6378514843 8447287804 2669665725 5067558074 3328222114 6997480992 3679095266
7723035108 1543450869 0659923574 6056699848 3716076746 7548415275 1995726963 8028406666 9346759480 9127972023 3445401568 2522725532
1694552698 2850467545 0582666679 1607539513 3020682543 9156668990 7317750218 7161752500 1719007676 3169499758 0549738598 5240661090
0997728324 5796817404 6771303866 3422757322 7197426540 4856247458 6210122019 9373992050 8173168442 1709024318 7223754580 3730348402
2597790757 0002133120 6040566400 9617913407 3320368260 5373421482 0759207570 6039643092 7389100660 3990780772 0589393831 0872494082
9493766287 8062899260 0224761029 2289257069 8480178527 8538183760 6213842880 0583915487 0359119158 0992251426 6982867334 3661419324
3273448404 6903443079 1595562780 6628207496 5421765530 5557100799 6196644907 8553956444 8839092547 3088067064 0750855438 0992248519
8461570976 7557676794 6424533172 8622263936 7937554904 7843037243 2758765658 4004386282 0558670280 2343937225 7391079717 3777956197
0726375440 5079901376 7408086919 3042793588 8398380961 2101825008 8390512187 7775244306 5670175011 4144853241 2968274812 9779440413
1531786063 5826826819 7144360048 8551917217 4152030009 2337959536 8604280741 7754077438 4187057463 1405188784 0579300617 0161081598
1046839400 4580506925 2928856195 7121352538 4170032274 7894253280 5116949543 4671439118 8412385641 9833844825 2778492706 8624396824
4544153875 7629726380 7003359549 8400891613 5071548163 4681971108 6615519015 6879197536 6906443556 1769879746 5676000325 3484630333
9582856997 2586363867 9179097742 5187820284 1891806261 8825410851 4671272988 1927869464 6209246153 8740736062 9976894156 9378765828
0929278776 4758755692 0334436137 9409594818 7898901673 3806245247 8392983900 6339982231 0836766874 9648037648 8163864933 2561015432
1038656410 0015072702 9951510388 8395034225 5356796067 9385576086 6058907138 1411887394 6560071238 2602303016 3441414947 2831067290
3058164329 1134572416 2603198046 8980257034 2409898609 9448088672 4106622116 0016713451 6232201273 1887440162 5340202243 0679976216
2699385365 8805932581 5792267039 8192560318 7514857779 2381741933 2711142002 5295433041 0790425638 3877548986 3329477270 5739109106
6397127994 3374463492 9096558678 2383764163 0764208573 5732862067 4712943103 6507332406 4787281063 9161542967 1934221388 2479326113
9141510120 6939178627 9970522085 2359786365 1463890464 8596938225 2501389808 8889951445 3543366890 6312178029 1612183169 7004020143
1268193999 8855590082 3977461330 4022328229 2192926811 3254803988 3458529504 3538394937 9407800834 2699889465 7353134284 2352604299
4711322324 9407155526 4475819954 7603440262 8307579115 2768890840 3332694257 7261987869 6453401334 2289385387 0880189222 4203712840
1949634095 2208208046 9143303493 0208953872 7100060736 9031310535 9679258030 0219168646 4629156732 1735420462 1765817159 9734688636
0489558914 6255768917 6988053128 8104551155 8779966874 1501796623 8491214989 6078412044 4537618939 7082364273 3075741794 4007728972
6581560958 3962706995 9576609368 4666605070 8889458468 7056610324 6549632496 9565641890 5920267892 2774161927 9818058791 7193546203
0564404995 5017218148 3005357653 4802210704 0414978603 8303734997 5760036852 5199409194 7600762233 3127755340 4869286981 0951538882
4714145597 4824197665 8048621154 4781022060 9891478028 5415513286 3649160416 6196380146 9041248750 5833555549 9131509416 1987429893
5209429728 8786231692 9835817894 8616294731 2101411348 2240418985 1350677733 8076678824 3223137419 0771050393 5893741695 1522625511
0387852732 5589924620 2463992289 9676049989 9361539534 4436847134 5324691104 6004956355 6972867972 2551148215 4027796976 4100205187
7853083338 4656546506 4147802343 5543703620 1288860208 2097667463 2834848111 9850881451 6127629067 4697181501 4852694458 7019720274
6986466944 8191082983 0579665223 5639939937 7470904670 6588659673 5000375820 3652785146 6082420420 5989637261 1839702565 3094783874
4778601361 9261275283 7475731804 0450908475 2855537107 4661019674 5930606064 8905747311 1069430685 7692886192 0946796030 9656111916
8565636178 5827399226 1122018531 3177156271 5943155708 3927352620 1546931057 8117588495 2258452893 6441636843 2485226778 9745031355
9729496259 0972586006 4224464682 8469110215 3636889005 2452847760 0452325742 7012058040 5741259039 1137024977 4892212034 0896125638
3057450610 9094576131 1635556372 3316096318 6185503601 3881252490 9556234366 0470839669 2226042439 8721120058 5522944124 6551784423
0690447551 7820798233 9105328582 7575180708 9477817877 3612154368 2819532455 3363374324 0739634223 1095615547 4206361954 0069012406
0121091857 4296394911 4768969975 0001995603 7938139759 2485841691 4959013234 5080597498 0393897360 0996614241 5529057573 1596330911
4323842109 9106487888 9098396239 6581582548 8383893031 8166049955 1020622178 4119170423 9679085446 6606943724 7345840146 2255913641
6352220939 0852695148 1591877095 9584078735 9507341430 3818911891 0864768198 4380468592 9030880738 1404177285 8403018502 3082686988
2383399363 2770798623 8834683835 1682906086 7665801777 6964197201 4979476022 1672976634 0428375883 5729504036 0131228789 9199452905
9149921275 8857505722 6316227338 6053355636 2193761119 2519522975 0569910611 0087641384 9097210121 5817055351 1489557026 5438293816
5699973717 8867209697 6784525445 0302046978 7996191291 8675827480 6039410787 3725265650 1147466483 7012074053 9637593311 1238703961
8306205130 3129727667 3315608991 4820790484 2140866443 5788979033 4576810283 0449238308 8649750111 8256619978 1606746012 9759166681
0637086010 3919549607 2630836634 6976395495 8166128284 7162999653 3299187898 9443895704 8831641016 7400040843 0845585070 4257823537
2201799745 0637678784 6015499193 6815960612 7309722714 2994666944 5769812682 1863556993 7497247380 4539858483 9761979506 4142082404
8883761778 8449858983 7649237616 9120629819 5129918568 2094479298 5103414696 0275280506 3374763647 4769488892 6720156330 2520670193
9541775633 7631005803 8146600217 5700696090 3334543510 2419377864 9423850468 1484224944 8334365336 6364845413 7021923280 6195803570
8759230170 7530891960 3136129863 1136242871 3507887120 6941583041 6695285030 0818903450 8637230221 3642474270 7355639721 8472668981
8600294854 6130984256 4393024589 9946305328 9399380016 0492898863 7349065378 0230221870 0449526155 6361521820 6006220481 0255571825
1169670216 0250130781 9654133270 7244944449 2173942182 1661554649 1144484300 4328245968 0112330992 9340785657 6694228949 0692749492
9076045179 5837958010 7181039399 4290665567 5741199429 4034750410 2078024331 1587966741 4137423494 9762095819 7007251919 9544039439
1035138339 0080135241 3504455649 3690651812 5576654432 2520810351 8153017971 9701503808 7034617108 6994560577 3005992089 4354771247
8984331830 1357943685 0286225359 4344734509 3929596596 3679278183 4846785904 7291946512 6334673736 5089417378 0200118304 8558809023
2707065869 0145899735 5392399886 4159485812 3461509148 7018473096 2357145969 7960093622 4252162608 8783434385 0156687645 9935767277
4309950181 7177129660 4503325324 6589910046 8248212061 3368563305 6919946692 0991435553 2945894753 9638830798 3021812182 7819389601
7573065789 4794218547 8535560420 4248416080 1953546817 5937699475 5776217305 1441262581 6973662980 0784056731 9554043028 7046015860
8271010226 4839353486 2545418255 9840777921 0024441338 2181317426 1962011650 4399922957 4854026661 1734556360 31
\end{Verbatim}

\end{document}
